% =====================================================================
\section{Was ist \LaTeX?}
% =====================================================================

\begin{frame}{Markup statt WYSIWYG}
  \begin{columns}[T]
    \begin{column}{0.47\textwidth}
      \begin{kasten}{Word, LibreOffice, \dots}
        \begin{itemize}
          \item \emph{What You See Is What You Get}
          \item Formatierung direkt beim Schreiben
          \item Layout liegt in eurer Verantwortung
        \end{itemize}
      \end{kasten}
    \end{column}
    \begin{column}{0.47\textwidth}
      \begin{kasten}{\LaTeX}
        \begin{itemize}
          \item Ihr beschreibt die \alert{Struktur}
          \item \LaTeX{} übernimmt das \alert{Layout}
          \item Inhalt und Form sind getrennt
        \end{itemize}
      \end{kasten}
    \end{column}
  \end{columns}
  \vspace{0.6em}
  \begin{merke}
    Ihr schreibt \texttt{\textbackslash section\{Einleitung\}}, nicht
    \glqq Arial, 14\,pt, fett, Abstand 12\,pt\grqq.
  \end{merke}
\end{frame}

\begin{frame}{Vom Quelltext zum PDF}
  \centering
  \vspace{1em}
  \begin{tikzpicture}[
      box/.style={rounded corners=5pt, fill=white, draw=hell, line width=1pt,
                  minimum width=3.2cm, minimum height=1.6cm, align=center},
      pfeil/.style={-{Stealth[length=3.5mm]}, line width=1.6pt, akzent}]
    \node[box] (tex) {\textbf{\texttt{bericht.tex}}\\[3pt]\scriptsize\color{grau}Quelltext};
    \node[box, right=1.6cm of tex, fill=dunkel, draw=dunkel, text=white] (comp)
      {\textbf{\texttt{pdflatex}}\\[3pt]\scriptsize\color{white!70!dunkel}Compiler};
    \node[box, right=1.6cm of comp] (pdf) {\textbf{\texttt{bericht.pdf}}\\[3pt]\scriptsize\color{grau}Ergebnis};
    \draw[pfeil] (tex) -- (comp);
    \draw[pfeil] (comp) -- (pdf);
    \node[below=0.5cm of comp, text=grau, font=\scriptsize]
      {Nebenprodukte: \texttt{.aux}, \texttt{.log}, \texttt{.toc}, \dots};
  \end{tikzpicture}
  \vspace{1.5em}
  \begin{itemize}
    \item Ähnlich wie beim Programmieren: schreiben, kompilieren, Ergebnis ansehen
    \item Fehler stehen in der \texttt{.log}-Datei bzw.\ im Editor
  \end{itemize}
\end{frame}

\begin{frame}{Geschichte und Relevanz}
  \begin{columns}[T]
    \begin{column}{0.47\textwidth}
      \textbf{Geschichte}\par\vspace{0.4em}
      \begin{itemize}
        \item \alert{1978:} Donald Knuth entwickelt \TeX, weil ihm der Satz seiner Bücher nicht gefällt
        \item \alert{1984:} Leslie Lamport baut \LaTeX{} als Makropaket darauf
        \item \alert{Heute:} \LaTeXe{} plus tausende Pakete auf CTAN
      \end{itemize}
    \end{column}
    \begin{column}{0.47\textwidth}
      \textbf{Warum in der Physik?}\par\vspace{0.4em}
      \begin{itemize}
        \item Formelsatz in hoher Qualität
        \item Standard für Paper, arXiv und Abschlussarbeiten
        \item Skripte und Übungsblätter sind damit gesetzt
        \item Kostenlos, plattformunabhängig, reine Textdateien
      \end{itemize}
    \end{column}
  \end{columns}
  \vspace{0.3em}
  \begin{tcolorbox}[enhanced, colback=white, colframe=hell, boxrule=0.8pt, arc=4pt, halign=center]
    $\displaystyle i\hbar\frac{\partial}{\partial t}\Psi(\mathbf{r},t)
      = \left[-\frac{\hbar^2}{2m}\nabla^2 + V(\mathbf{r},t)\right]\Psi(\mathbf{r},t)$
  \end{tcolorbox}
\end{frame}
